\section{Methodology}
\label{sec:method}

\subsection{Overview}

Our experimental pipeline has four stages: data generation, representation, training, and evaluation.
We deliberately fix all hyperparameters across cipher families so that any performance difference is attributable to cipher structure, not tuning.

\subsection{Data Generation}

For each cipher at $r$ rounds, we generate a balanced dataset of $N = 500{,}000$ ciphertext pairs.
Positive samples (label 1) encrypt $(P, P \oplus \Delta P)$ under a uniformly random key $K$; negative samples (label 0) encrypt two independently random plaintexts under the same $K$, following Gohr's protocol~\citep{gohr2019}.
Each experiment uses 3--5 random seeds; the key changes with each seed.

\subsection{Input Representation}
\label{sec:representation}

We expand each ciphertext word into individual bits and concatenate:
\begin{equation}
  \mathbf{x} = \operatorname{bits}(C_L) \| \operatorname{bits}(C_R) \| \operatorname{bits}(C'_L) \| \operatorname{bits}(C'_R) \| \operatorname{bits}(C_L \oplus C'_L) \| \operatorname{bits}(C_R \oplus C'_R),
\end{equation}
yielding a 96-dimensional input for 32-bit block ciphers and 192-dimensional for \present (64-bit).
We denote this the \emph{R2\_xor\_diff} representation.

\paragraph{Justification.}
Among six representations tested on \speck at 5 rounds (E02; \Cref{sec:app-ablations}), R2\_xor\_diff achieves 88.7\% accuracy, tied with the manually engineered R8\_statistical (88.6\%) and significantly above the word-level R5 baseline (59.8\%).
On \present at 5 rounds, R2\_xor\_diff (76.1\%) again ties with R8\_statistical (76.1\%) and dominates all others.
On \simon at 7 rounds, R4\_bit\_sliced (85.6\%) outperforms R2\_xor\_diff (81.3\%) by 4.3\pp, suggesting that cipher-specific representation tuning could improve absolute accuracy.
We retain R2\_xor\_diff across all families to preserve the controlled comparison that is our primary contribution.

\subsection{Architecture}
\label{sec:architecture}

Our primary architecture is a feedforward MLP: four hidden layers of sizes $[256, 128, 64, 32]$ with batch normalization after each, followed by a single sigmoid output.
This is Gohr's simpler MLP variant~\citep{gohr2019}, deliberately chosen as a \emph{lower bound} on neural distinguisher performance.

We also evaluate six additional architectures---CNN, Residual CNN, Gohr-MLP (larger), LSTM, GRU, and Siamese---on \speck at 5 rounds, as well as Gohr's depth-10 ResNet on all three ciphers (\Cref{sec:exp-architecture}).

\subsection{Training Protocol}

All models use Adam ($\text{lr} = 10^{-3}$), binary cross-entropy loss, batch size 5000, for up to 30 epochs with early stopping (patience 5) on validation accuracy.
Training data is regenerated per seed.
All experiments run on a single NVIDIA Tesla V100-PCIE-32GB.

\subsection{Classical Baseline}

As a classical reference, we implement a best-bit-bias distinguisher.
For each of the $2b$ output difference bits, we compute $\epsilon_i = |p_i - 0.5|$ over $10^5$ samples across 5 keys.
The distinguisher classifies using the bit with maximum bias.
This is a single-bit distinguisher, intentionally weaker than DDT-based methods~\citep{gohr2022assessment}; we report it as a lower bound on classical performance and cite Gohr's published numbers for the stronger ResNet comparison.
